\chapter*{不定积分每日一题}

\refer{《积分的方法与技巧》金玉明\ 顾新身\ 毛瑞庭\ /\ 编著}

\exercise{求积分$\int{\frac{3x+5}{x^3-x^2-x+1}\d x}$}

\teacher{夏明宇}

\postdate{2022年11月29日}

\result{\frac{1}{2}\ln{\abs{\frac{x+1}{x-1}}}-\frac{4}{x-1}+C}

\analysis{经典的有理分式积分，分解后的待定系数确定常用赋值法减少计算.}

\solution{
    被积函数的分母可分解为（赋值找到因子+大除法得到）
    \formula{x^3-x^2-x+1=\qty(x+1)\qty(x-1)^2}
    因此，被积函数可写成（不明白的同学可以看一下教材有理积分部分哦）\formula{\frac{3x+5}{x^3-x^2-x+1}=\frac{A}{x+1}+\frac{B}{x-1}+\frac{C}{\qty(x-1)^2}}
    两边同时乘以$\qty(x+1)\qty(x-1)^2$，得
    \formula{3x+5=A\qty(x-1)^2+B\qty(x+1)\qty(x-1)+C\qty(x+1)}
    这时，我们只需要确定这些系数，这道积分就可以迎刃而解了.


    比较朴素的办法是右侧展开后和左侧进行比较来确定三个常数，但这样做比较麻烦. 我们考虑这个求系数的本质：恒等式.左右两边是恒成立的，因此我们可以利用这一性质 进行操作.比较常见和好用的是赋值.如本题中，令$x=1$，使含$A$、$B$的两项为零，得$C=4$；再令$x=-1$，使含$B$、$C$的两项为零，得$A=\frac{1}{2}$；最后发现左侧没有$x^2$项，右侧$x^2$项的系数为$A+B$所以$B=-A=-\frac{1}{2}$，那么剩下的就好做了：
    \formula{
        \begin{aligned}
            \int{\frac{3x+5}{x^3-x^2-x+1}\d x}&=\int{\qty[\frac{1}{2}\frac{1}{x+1}-\frac{1}{2}\frac{1}{x-1}+4\frac{1}{\qty(x-1)^2}]\d x}\\
            &=\frac{1}{2}\ln{\abs{x+1}}-\frac{1}{2}\ln{\abs{x-1}}-\frac{4}{x-1}+C\\
            &=\frac{1}{2}\int{\frac{\d x}{x+1}-\frac{1}{2}\int{\frac{\d x}{x-1}}+4\int{\frac{\d x}{\qty(x-1)^2}}}\\
            &=\frac{1}{2}\ln{\abs{\frac{x+1}{x-1}}}-\frac{4}{x-1}+C
        \end{aligned}
    }
    其实不只是赋值，因为本质是恒等式，因此我们也可以对两边同时求导，同时取极限 等，而这些操作的最终目的就是化简计算，哪个简单哪个来，也不必拘泥于一种形式.
}